% Relative field accesses: offsets, compass names, and index dimensions.
\begin{tikzpicture}[font=\footnotesize]
  % ---- (a) scalar field u: the 3x3 neighbourhood -------------------------------
  \begin{scope}[x=1.55cm,y=1.15cm]
    \foreach \dx/\dy/\nm in {-1/1/NW,0/1/N,1/1/NE,-1/0/W,1/0/E,-1/-1/SW,0/-1/S,1/-1/SE} {
      \fill[psorange!25] (\dx+1,\dy+1) rectangle ++(1,1);
      \node[align=center,font=\scriptsize] at (\dx+1.5,\dy+1.5)
        {\texttt{u[\dx,\dy]}\\[-1pt]\textcolor{psgray}{\nm}};
    }
    \fill[psblue!55] (1,1) rectangle ++(1,1);
    \node[align=center,font=\scriptsize,text=white] at (1.5,1.5)
      {\texttt{u[0,0]}\\[-1pt]\texttt{u.center}};
    \draw[gridline] (0,0) grid (3,3);
    \draw[axisarrow] (0,-0.25) -- (1,-0.25) node[right,celltext] {$x$ / first offset};
    \draw[axisarrow] (-0.12,0) -- (-0.12,1) node[above,celltext] {$y$};
    \node[font=\small] at (1.5,3.35) {(a) \texttt{u = ps.fields("u: double[2D]")}};
  \end{scope}

  % ---- (b) vector field v(2): one value per index per cell -----------------------
  \begin{scope}[shift={(6.6cm,0.35cm)},x=1.55cm,y=0.75cm]
    \foreach \i in {0,1,2} {
      \draw[gridline] (\i,0) rectangle ++(1,2);
      \draw[gridline] (\i,1) -- ++(1,0);
    }
    \fill[psblue!25] (1,1) rectangle ++(1,1);
    \fill[nbcell] (2,1) rectangle ++(1,1);
    \foreach \i in {0,1,2} {
      \node[celltext] at (\i+0.5,0.5) {\texttt{(0)}};
      \node[celltext] at (\i+0.5,1.5) {\texttt{(1)}};
      \draw[domainline] (\i,0) rectangle ++(1,2);
    }
    \node[celltext] at (0.5,-0.3) {W};
    \node[celltext] at (1.5,-0.3) {C};
    \node[celltext] at (2.5,-0.3) {E};
    \draw[psgray] (1.5,2) -- (1.5,2.6) node[above,celltext] {\texttt{v(1)} = \texttt{v.center(1)}};
    \draw[psgray] (2.5,2) -- (2.9,3.3) node[above,celltext] {\texttt{v[1,0](1)}};
    \draw[decorate,decoration={brace,amplitude=3pt}] (-0.08,0) -- (-0.08,2)
      node[midway,left=4pt,align=right,celltext] {index\\dimension\\(size 2)};
    \node[font=\small] at (1.5,4.67) {(b) \texttt{v = ps.fields("v(2): double[2D]")}};
  \end{scope}
\end{tikzpicture}
